\chapter{SYSTEM DESIGN}

\ChapterSection{Design}

Four components move the system from the requirements in Chapter~3 to a deployable diagnosis
pipeline: a YOLO detector for plant-part localization, a fine-tuned VLM for disease
classification, a backend that orchestrates inference and advisory generation, and a mobile
application that delivers results to the farmer. Each component is presented below with its
datasets and configuration, followed by the algorithmic procedures that run across them.

\ChapterSubsection{YOLO Object-Detection Dataset}

The YOLO component localizes leaves, panicles, grain clusters, fruits, ears, and vegetable
heads within crop images. Disease classification is left to the VLM --- the detector's job is
to tell the downstream pipeline which plant parts are present and where they sit in the frame.

\begin{figure}[htbp]
\centering
\includegraphics[width=1\textwidth,height=0.5\textheight,keepaspectratio]{data/flow_digram_data/build/yolo/06_system_architecture_overview.pdf}
\caption{YOLO system architecture overview showing the data, training, validation, deployment, and export pipeline.}
\label{fig:yolo-system-architecture-overview}
\end{figure}
YOLOv8 nano fits the deployment target: it is small enough for a mobile TFLite export while
reaching usable mAP on the eleven-class taxonomy
\parencite{jiangYOLOSurvey2022,ultralyticsYolov8}.
Figure~\ref{fig:yolo-system-architecture-overview} outlines the detection pipeline from data
ingestion through export.

\subsubsection{Data Sources and Class Mapping}

Sources span Kaggle and Roboflow datasets covering tomato, rice, potato, maize, brassica,
cauliflower, and cabbage imagery, merged into a single class taxonomy. From 34,992
candidate images, curation removed 4,011 images and 13,218 annotations flagged for
out-of-bounds coordinates, sub-pixel area, or collapse after clamping. The class distribution
is heavily skewed: rice panicle holds 115,046 of 147,467 annotations, while maize ear
(106) and rice grain cluster (473) are the least represented. Aggregate mAP hides this
imbalance; per-class metrics must be read from individual class scores
(Figure~\ref{fig:yolo-data-pipeline-flow}).

\begin{figure}[htbp]
\centering
\includegraphics[width=1\textwidth,height=0.5\textheight,keepaspectratio]{data/flow_digram_data/build/yolo/01_data_pipeline_flow.pdf}
\caption{YOLO data pipeline flow from raw annotations to canonical export, splitting, and training structure.}
\label{fig:yolo-data-pipeline-flow}
\end{figure}

\subsubsection{Dataset Construction}

The final dataset contains 30,981 images and 147,467 annotations across eleven classes,
split into 27,454 training and 3,527 validation images. Annotations use COCO-style pixel
\texttt{xywh} format, converted to YOLO format for training. Per-class counts appear in
Table~\ref{tab:yolo-dataset-class-distribution}.

{\small
\renewcommand{\arraystretch}{1.3}
\begin{xltabular}{\textwidth}{@{}lr@{}}
\caption{YOLO dataset class distribution after curation}\label{tab:yolo-dataset-class-distribution}\\
\toprule
\textbf{Class} & \textbf{Object annotations} \\
\midrule
\endfirsthead
\multicolumn{2}{l}{\footnotesize\textit{Table \thetable{} \textemdash{} continued from previous page}}\\[3pt]
\toprule
\textbf{Class} & \textbf{Object annotations} \\
\midrule
\endhead
\endfoot
\bottomrule
\endlastfoot
Rice leaf & 1,657 \\
Rice panicle & 115,046 \\
Rice grain cluster & 473 \\
Maize leaf & 8,000 \\
Maize ear & 106 \\
Potato leaf & 1,199 \\
Tomato leaf & 8,792 \\
Tomato fruit & 6,808 \\
Brassica leaf & 3,834 \\
Cauliflower head & 525 \\
Cabbage head & 1,027 \\
\end{xltabular}
}

The detector runs on 640-pixel inputs and exports to FP32, FP16, and INT8 TFLite variants
for mobile deployment; training and threshold-selection details follow in Algorithm Details.

\ChapterSubsection{VLM Disease-Diagnosis Dataset}

Disease interpretation runs on Qwen2.5-VL-3B \parencite{bai2025qwen25vl}. The model takes a
crop image and a natural-language instruction, then returns a structured JSON diagnosis. At 3B
parameters with 4-bit LoRA fine-tuning, it fits on a single L4 GPU. Because the output is
generated text --- not a fixed-vocabulary prediction --- parseability is a first-class
evaluation metric alongside classification accuracy
(Figure~\ref{fig:vlm-system-architecture-overview}).

\begin{figure}[htbp]
\centering
\includegraphics[width=0.9\textwidth,height=0.45\textheight,keepaspectratio]{data/flow_digram_data/build/vlm/06_system_architecture_overview.pdf}
\caption{VLM system architecture overview separating dataset processing, training, inference, and serving infrastructure.}
\label{fig:vlm-system-architecture-overview}
\end{figure}

\subsubsection{Taxonomy and Label Normalization}

The taxonomy covers rice, maize, potato, tomato, cauliflower, and cabbage across eleven
disease classes: maize common rust, maize healthy, maize northern leaf blight, potato early
blight, potato healthy, potato late blight, tomato early blight, tomato healthy, tomato late
blight, tomato leaf mold, and tomato septoria leaf spot. Canonical labels follow the format
\texttt{crop\_\_disease} (e.g., \texttt{potato\_\_late\_blight}). Evaluation uses exact
canonical-label matching, so alias normalization at export is a correctness requirement, not a
cosmetic one.

\subsubsection{Dataset Pipeline}

The dataset pipeline is configured under \texttt{vlm\_model/data/configs} as project
\texttt{krishivaidya}, dataset \texttt{kv-core5-ms}. Splits reserve 15\% for a locked test
set and 15\% for validation. Per-crop record caps --- 800 for rice, maize, potato, and tomato;
600 for cauliflower and cabbage --- prevent any single crop from dominating training
(Figure~\ref{fig:vlm-data-pipeline-flow}).

\begin{figure}[htbp]
\centering
\includegraphics[width=1\textwidth,height=0.6\textheight,keepaspectratio]{data/flow_digram_data/build/vlm/01_data_pipeline_flow.pdf}
\caption{VLM data pipeline flow from raw sources through curation, canonicalization, export, and training data assembly.}
\label{fig:vlm-data-pipeline-flow}
\end{figure}

\subsubsection{Qwen2.5-VL SFT Export}

Each exported record pairs an image path with a human message (containing the image token and
instruction), an assistant answer, and metadata (split, crop, disease, canonical label,
template id, source). Three templates shape the assistant target: \texttt{disease\_json}
requests compact JSON with \texttt{crop}, \texttt{disease}, \texttt{canonical\_label}, and
\texttt{confidence\_label}; \texttt{label\_only} requests the canonical and confidence labels
alone; \texttt{diagnostic\_short} adds a brief visual-diagnosis sentence.

LoRA adapters target the attention and feed-forward projection layers, excluding visual
modules; fine-tuning configuration and prediction parsing follow in Algorithm Details.

\ChapterSubsection{Backend System Architecture}

A TypeScript and Express API sits between the mobile client and every backend concern: farm
records, image storage, authentication, VLM inference, and advisory generation, all behind a
single versioned interface. The codebase separates configuration, middleware, routing,
controllers, services, repositories, and data models into layers so that HTTP transport,
business rules, database persistence, and ML integrations do not leak into each other
\parencite{sculleyMLTechDebt2015}.

Figure~\ref{fig:backend-architecture} traces the request path: the mobile app sends JSON and
multipart requests to Express; controllers validate inputs, services enforce workflow rules,
repositories wrap Mongoose queries, and models define MongoDB documents.

\begin{figure}[htbp]
\centering
\includegraphics[width=0.9\textwidth,height=0.45\textheight,keepaspectratio]{data/flow_digram_data/build/backend/05_system_architecture_overview.pdf}
\caption{Backend system architecture overview showing the request stack and external service boundaries.}
\label{fig:backend-architecture}
\end{figure}

\subsubsection{Runtime Modes and Application Initialization}

Three runtime modes control which route groups load, set through the \texttt{BACKEND\_MODE}
environment variable: \texttt{diagnose} (AI routes only), \texttt{non\_ai} (database routes
only), and \texttt{complete} (all routes). A health endpoint at \texttt{/api/v1/health}
reports server state, uptime, environment, mode, and public base URL. Zod validates all
configuration before the server starts, preventing silent startup in an incomplete state.
Database-backed routes require JWT, Cloudinary, and SMTP credentials; the diagnosis stack
requires the ML service URL and its authentication configuration.

Dependencies are wired in \texttt{src/container.ts} --- repositories first, then services,
controllers, and route modules --- keeping the dependency graph in one place instead of
scattered across route files.

\ChapterSubsection{Backend API Design}

\subsubsection{Route Groups}

All routes sit under \texttt{/api/v1} (Table~\ref{tab:backend-route-groups}). In development,
Swagger UI at \texttt{/api/docs} exposes request schemas, security requirements, and response
documentation generated from route-level comments.

{\small
\renewcommand{\arraystretch}{1.3}
\begin{xltabular}{\textwidth}{@{} l L L @{}}
\caption{Implemented backend route groups}\label{tab:backend-route-groups}\\
\toprule
\textbf{Route group} & \textbf{Base path} & \textbf{Implemented responsibility} \\
\midrule
\endfirsthead
\multicolumn{3}{l}{\footnotesize\textit{Table \thetable{} \textemdash{} continued from previous page}}\\[3pt]
\toprule
\textbf{Route group} & \textbf{Base path} & \textbf{Implemented responsibility} \\
\midrule
\endhead
\endfoot
\bottomrule
\endlastfoot
Health & \texttt{/api/v1/health} & Server status, runtime mode, uptime, and exposure information \\
Diagnosis & \texttt{/api/v1/diagnose} & Image validation, VLM inference request, advisory generation, diagnosis response \\
Authentication & \texttt{/api/v1/auth} & Registration, login, token refresh, logout, OTP, password reset \\
Users & \texttt{/api/v1/users} & Profile retrieval, profile update, profile image, password change, push tokens \\
Scans & \texttt{/api/v1/scans} & Scan upload, scan history, statistics, grouping, status, resolution, follow-up linkage \\
Crops & \texttt{/api/v1/crops} & Crop-plot creation, filtering, archive state, growth stage, health status \\
Tasks & \texttt{/api/v1/tasks} & Crop-task creation, batch creation, pending/overdue/upcoming views, completion state \\
Notes & \texttt{/api/v1/notes} & Crop notes, linked scan notes, note images, note-type counts \\
Advisories & \texttt{/api/v1/advisories} & Advisory creation, filtering, active advisories, dismiss/undismiss actions \\
\end{xltabular}
}

\subsubsection{Controller and Validation Pattern}

Controllers take typed service dependencies and expose Express request handlers. Validation
uses \texttt{express-validator} for MongoDB identifier checks, bounded pagination, enumerated
types (task categories, advisory severities), and text-field length limits. Successful
responses pass through a shared formatter; failures reach a global error middleware.

Most database-backed routes require authentication. The auth middleware verifies a Bearer token
against the JWT secret, rejects expired or malformed tokens, blocks pending-verification and
suspended accounts, and attaches the decoded user context to the request. Service methods then
enforce ownership through repository queries, so access control does not repeat in every
controller.

\ChapterSubsection{Backend Data Model}

\subsubsection{MongoDB Collections and Domain Entities}

Mongoose models back nine domain entities: users, OTP records, refresh tokens, scans, scan
resolutions, crop plots, crop tasks, crop notes, and advisories
(Table~\ref{tab:backend-data-models}). Indexes cover user ownership, status filtering, date
ordering, and TTL cleanup for expiring documents.

{\small
\renewcommand{\arraystretch}{1.3}
\begin{xltabular}{\textwidth}{@{} l L L @{}}
\caption{Implemented backend data models}\label{tab:backend-data-models}\\
\toprule
\textbf{Model} & \textbf{Primary purpose} & \textbf{Important implementation details} \\
\midrule
\endfirsthead
\multicolumn{3}{l}{\footnotesize\textit{Table \thetable{} \textemdash{} continued from previous page}}\\[3pt]
\toprule
\textbf{Model} & \textbf{Primary purpose} & \textbf{Important implementation details} \\
\midrule
\endhead
\endfoot
\bottomrule
\endlastfoot
\texttt{User} & User account and preferences & Email, password hash, name, language, region, accessibility, location, status, push tokens \\
\texttt{Otp} & OTP verification state & Purpose-specific OTP hash, attempt count, expiry, unique email-purpose index, TTL expiry \\
\texttt{RefreshToken} & Session refresh state & Token hash, lookup ID, device context, revocation state, expiry, TTL cleanup \\
\texttt{Scan} & Disease scan history & Image URI, crop metadata, status, result object, client timestamp, offline flag, soft delete \\
\texttt{ScanResolution} & Outcome tracking for scans & Resolution type, treatment used, notes, follow-up scan references \\
\texttt{CropPlot} & Farmer crop plot record & Crop type, variety, region, district, area, planting date, stage, health status, archive state \\
\texttt{CropTask} & Farm task and reminder record & Plot link, category, due date, priority, completion state, auto-generated flag \\
\texttt{CropNote} & Crop observation and treatment note & Plot link, optional scan link, note type, title/body, image attachments, soft delete \\
\texttt{Advisory} & Crop advisory notification & Plot link, advisory type, severity, Nepali/English fields, dismiss state, expiry TTL \\
\end{xltabular}
}

\subsubsection{Repository Layer}

Repositories encapsulate database access for user-owned resources --- scans, crop plots,
tasks, notes, and advisories. Each method accepts a user or plot identifier and applies
ownership filters before returning results, so access control lives in the data layer rather
than in individual controllers (Figure~\ref{fig:backend-data-model-persistence-flow}).

\begin{figure}[!htbp]
\centering
\includegraphics[width=0.9\textwidth,height=0.40\textheight,keepaspectratio]{data/flow_digram_data/build/backend/03_data_models_persistence_flow.pdf}
\caption{Backend data model and persistence flow showing the relationship between controllers, services, repositories, and MongoDB collections.}
\label{fig:backend-data-model-persistence-flow}
\end{figure}

\ChapterSubsection{Backend Diagnosis and Advisory Pipeline}

\subsubsection{Image Validation}

Image validation runs before model inference: malformed files, unsupported formats,
out-of-range dimensions, and invalid crop hints are rejected before they reach the VLM or
advisory services.

\subsubsection{VLM Inference Bridge}

\begin{figure}[htbp]
\centering
\includegraphics[width=0.9\textwidth,height=0.4\textheight,keepaspectratio]{data/flow_digram_data/build/vlm/04_api_gateway_flow.pdf}
\caption{VLM API gateway flow showing the backend request path, readiness check, authentication, inference call, and response parsing boundary.}
\label{fig:vlm-api-gateway-flow}
\end{figure}
The VLM inference gateway is a FastAPI service \parencite{fastapiDocs} that exposes
\texttt{/health}, \texttt{/ready}, and \texttt{/v1/inference}. The Express backend holds the
gateway credentials and checks readiness before forwarding any inference request
(Figures~\ref{fig:vlm-api-gateway-flow} and~\ref{fig:backend-diagnosis-pipeline-flow}).

\begin{figure}[htbp]
\centering
\includegraphics[width=0.9\textwidth,height=0.4\textheight,keepaspectratio]{data/flow_digram_data/build/backend/02_diagnosis_pipeline_flow.pdf}
\caption{Backend diagnosis pipeline flow showing image validation, VLM inference bridging, advisory generation, and response assembly.}
\label{fig:backend-diagnosis-pipeline-flow}
\end{figure}

\subsubsection{Google AI Advisory Bridge}

The advisory bridge demands a strict JSON response covering disease, confidence,
out-of-distribution state, symptoms, treatment, prevention, safety, source, and
alternative-diagnosis fields. Zod schema validation gates the response before it reaches the
mobile application.

\ChapterSubsection{Backend Security, Uploads, and Reliability Controls}

\subsubsection{Upload Handling and Cloudinary Storage}

Multer memory storage handles incoming files, filtered by MIME type and extension against
JPEG, PNG, and WEBP. Scan images, profile pictures, and note images upload to Cloudinary with
secure URLs and automatic quality and format transformations. When a profile picture or note
image is replaced or removed, the backend deletes the previous Cloudinary asset.

\subsubsection{Request Safety and Error Handling}

Rate limits cap general routes at 100 requests per 15 minutes, authentication routes at 10,
and OTP routes at 5 per hour. A global error handler folds application errors, Multer upload
errors, Mongoose validation errors, duplicate-key errors, and JWT errors into a uniform JSON
shape. Production surfaces unknown errors as a generic message; development mode exposes full
diagnostic detail.

\subsubsection{Authentication and Session Handling}

Login returns a short-lived access token and stores a hashed refresh token in an HTTP-only
cookie. Refresh tokens carry lookup IDs, expiry timestamps, and revocation state, supporting
both single-session and all-session logout. OTP records store hashed values with expiry and
attempt counts under a unique email-purpose index, so each purpose --- email verification,
password reset, phone verification --- has its own expiration window.

\ChapterSubsection{Mobile Application Architecture}

\subsubsection{Purpose and Implementation Scope}

The mobile application is built with React Native and Expo
\parencite{reactNativeDocs,expoDocs} and handles onboarding, crop-plot management,
offline-aware scan queuing, diagnosis submission, result presentation, and bilingual record
keeping. A local-first design \parencite{kleppmannLocalFirst2019} writes user actions to
SQLite before synchronizing with the backend. The stack pairs Expo Router (navigation) with
TanStack Query (server-state mutations), i18next (bilingual content), VisionCamera (native
capture), and TFLite (on-device inference). The Android package is
\texttt{com.krishivaidya.app} under scheme \texttt{krishivaidya}
(Figure~\ref{fig:mobile-app-architecture}).

\begin{figure}[htbp]
\centering
\includegraphics[width=0.9\textwidth,height=0.45\textheight,keepaspectratio]{data/flow_digram_data/build/mobile_app/05_system_architecture_overview.pdf}
\caption{Mobile application architecture overview showing presentation, state, native modules, and local storage boundaries.}
\label{fig:mobile-app-architecture}
\end{figure}

\subsubsection{Application Shell and Navigation Structure}

The root layout initializes fonts, the splash screen, a QueryClient with offline-tolerant
defaults, i18n, and the repository bootstrap before rendering the route tree. Route protection
redirects first-time users through onboarding before they reach the main tabs.

Five bottom tabs --- Home, Crops, Scan, History, and Settings --- anchor the daily interface
(Table~\ref{tab:mobile-route-groups}). The scan stack (capture, preview, image editor,
analysis, final analysis) runs independently of crop management routes, so scan results and
reminders can cross-reference crop records without coupling the two flows.

{\small
\renewcommand{\arraystretch}{1.3}
\begin{xltabular}{\textwidth}{@{} l L L @{}}
\caption{Implemented mobile route and workflow groups}\label{tab:mobile-route-groups}\\
\toprule
\textbf{Route group} & \textbf{Implemented screens} & \textbf{Primary purpose} \\
\midrule
\endfirsthead
\multicolumn{3}{l}{\footnotesize\textit{Table \thetable{} \textemdash{} continued from previous page}}\\[3pt]
\toprule
\textbf{Route group} & \textbf{Implemented screens} & \textbf{Primary purpose} \\
\midrule
\endhead
\endfoot
\bottomrule
\endlastfoot
Onboarding & Welcome, language, permission, crop, demo & First-run setup, language selection, permission preparation, crop context \\
Main tabs & Home, Crops, Scan, History, Settings & Primary application navigation for daily use \\
Scan stack & Capture, preview, image editor, analysis, final analysis & Image acquisition, review, detection, cloud diagnosis, report viewing \\
Crop routes & Add/edit crop, crop detail, crop weather, tasks, notes & Farm plot management and agronomic record keeping \\
Settings routes & Profile, option select, notifications, accessibility, export, model precision & Preference management, data export, accessibility, model deployment controls \\
Achievements & Achievements view and toast components & Lightweight engagement tracking based on completed app actions \\
\end{xltabular}
}

\subsubsection{Provider, Repository, and Dependency Structure}

Two context providers supply dependencies to the component tree.
\texttt{ServicesProvider} handles camera, network, notification, haptics, voice, and location;
\texttt{RepositoryProvider} handles scans, crops, tasks, notes, and the data catalog. At
startup, the repository bootstrap initializes SQLite, runs the V3-to-V4 migration, performs a
one-time MMKV-to-SQLite migration, constructs datasources, creates repository instances, and
starts the offline sync manager. Read-only agricultural catalog data --- varieties, growth
stages, fertilizer recommendations, disease entries, and seasonal tips --- ships as JSON
fixtures, separate from user records in SQLite.

\subsubsection{Local Persistence and Offline-Aware Design}

The SQLite schema holds scans, offline queue entries, reminders, sync queue entries, crop
plots, tasks, notes, advisories, custom task categories, settings, achievements, achievement
statistics, and notifications. Mutable tables carry \texttt{sync\_status},
\texttt{updated\_at}, and \texttt{deleted\_at} columns so that writes complete locally before
backend synchronization runs (Table~\ref{tab:mobile-local-persistence}).

{\small
\renewcommand{\arraystretch}{1.3}
\begin{xltabular}{\textwidth}{@{} l L L @{}}
\caption{Implemented mobile local persistence groups}\label{tab:mobile-local-persistence}\\
\toprule
\textbf{Persistence group} & \textbf{Tables or stores} & \textbf{Supported manuscript interpretation} \\
\midrule
\endfirsthead
\multicolumn{3}{l}{\footnotesize\textit{Table \thetable{} \textemdash{} continued from previous page}}\\[3pt]
\toprule
\textbf{Persistence group} & \textbf{Tables or stores} & \textbf{Supported manuscript interpretation} \\
\midrule
\endhead
\endfoot
\bottomrule
\endlastfoot
Scan records & \texttt{scans}, \texttt{offline\_queue} & Local scan history, queued uploads, offline diagnosis fallback \\
Synchronization & \texttt{sync\_queue}, sync columns on mutable tables & Retryable operation queue and change tracking mechanism \\
Crop management & \texttt{crop\_plots}, \texttt{crop\_tasks}, \texttt{crop\_notes}, \texttt{crop\_advisories} & Farmer plot, task, note, and advisory records \\
Personalization & \texttt{user\_settings} & Language, theme, accessibility, onboarding, model precision, selected crops \\
Engagement and alerts & \texttt{user\_achievements}, \texttt{user\_achievement\_stats}, \texttt{notifications}, \texttt{reminders} & App engagement state, reminder scheduling, notification history \\
Catalog data & JSON fixture files through catalog repository & Read-only crop, growth-stage, fertilizer, disease, and seasonal reference data \\
\end{xltabular}
}

The offline sync manager monitors connectivity, resets interrupted items on startup, and
syncs every five minutes in batches of ten. Processing halts when connectivity drops. The
current table-push logic marks local rows as synced, but not all entity types have matching
backend endpoints yet --- so offline sync is working infrastructure with partial backend
coverage, not yet a fully bidirectional system
(Figure~\ref{fig:mobile-offline-sync-data-flow}).

\begin{figure}[htbp]
\centering
\includegraphics[width=0.9\textwidth,height=0.4\textheight,keepaspectratio]{data/flow_digram_data/build/mobile_app/03_offline_sync_data_flow.pdf}
\caption{Mobile offline synchronization data flow showing local writes, queue management, connectivity checks, and retry behavior.}
\label{fig:mobile-offline-sync-data-flow}
\end{figure}

\subsubsection{Scan Workflow and Diagnosis Flow}

The scan workflow moves through capture or gallery selection, preview, optional editing,
diagnosis submission, result display, and local history update
(Figure~\ref{fig:mobile-scan-pipeline-flow}). State transitions, network branching, timeout
handling, and retry logic are covered in the Algorithm Details section.

\begin{figure}[htbp]
\centering
\includegraphics[width=0.9\textwidth,height=0.6\textheight,keepaspectratio]{data/flow_digram_data/build/mobile_app/02_scan_pipeline_flow.pdf}
\caption{Mobile scan pipeline flow showing capture, preview, editing, diagnosis submission, and retry-aware request handling.}
\label{fig:mobile-scan-pipeline-flow}
\end{figure}

\subsubsection{On-Device Model Asset Support}

A singleton TFLite service manages three precision variants:
\texttt{krishi\_bouncer\_int8.tflite}, \texttt{krishi\_bouncer\_fp16.tflite}, and
\texttt{krishi\_bouncer\_fp32.tflite}, with runtime precision selection and delegate fallback
for Android inference (Figure~\ref{fig:mobile-ondevice-ml-inference-flow}).

\begin{figure}[htbp]
\centering
\includegraphics[width=0.9\textwidth,height=0.45\textheight,keepaspectratio]{data/flow_digram_data/build/mobile_app/04_ondevice_ml_inference_flow.pdf}
\caption{Mobile on-device ML inference flow showing model selection, load fallback, and runtime precision switching.}
\label{fig:mobile-ondevice-ml-inference-flow}
\end{figure}

\subsubsection{Localization, Accessibility, and Farmer-Facing Output}

Language defaults to Nepali, resolved from the SQLite preference first, then device locale,
then fallback. Result screens support treatment timelines, organic-alternative toggling,
reminder creation, share-card generation, PDF export, voice output, and feedback submission.

\ChapterSubsection{Integrated Diagnosis Workflow}

The four subsystems cooperate in one end-to-end path: the mobile scan pipeline captures and
routes the image, the backend validates it and calls the VLM gateway, the gateway returns a
prediction, and the advisory client produces a structured response that the mobile app presents
to the farmer. Figure~\ref{fig:integration-workflow} shows the combined flow.

\begin{figure}[htbp]
\centering
\includegraphics[width=0.9\textwidth,height=0.4\textheight,keepaspectratio]{data/flow_digram_data/build/00_all_flows.pdf}
\caption{Combined flow-diagram overview covering the subsystem workflows used in the manuscript.}
\label{fig:integration-workflow}
\end{figure}

Deployment splits runtime responsibilities across five nodes: the Android device, the Express
backend, MongoDB and Cloudinary storage, the Google AI advisory service, and the FastAPI VLM
gateway. Model-service credentials stay behind the backend boundary; YOLO TFLite variants ship
inside the mobile APK for on-device plant-part detection. The step-by-step online diagnosis
procedure is listed in Table~\ref{tab:integration-control-flow} under Algorithm Details.

\ChapterSection{Algorithm Details}

The procedures below operate across the components, data stores, and boundaries defined above.
Configuration values appear where they shape the procedure; measured outcomes are reported in
Chapter~5.

\ChapterSubsection{YOLO Training, Validation, and Deployment Procedure}

\subsubsection{Training Setup}

Training begins from the curated plant-part dataset. Image records and bounding-box
annotations are exported to YOLO format, then trained on \texttt{yolov8n.pt} with 640-pixel
inputs for 100 epochs. Mosaic augmentation runs at probability 1.0 and shuts off during the
final 10 epochs to reduce late-stage noise; mixup and copy-paste are disabled
(Table~\ref{tab:yolo-training-configuration}).

{\small
\renewcommand{\arraystretch}{1.3}
\begin{xltabular}{\textwidth}{@{}ll@{}}
\caption{YOLO training configuration used for the Krishi Vaidya detector}\label{tab:yolo-training-configuration}\\
\toprule
\textbf{Parameter} & \textbf{Value} \\
\midrule
\endfirsthead
\multicolumn{2}{l}{\footnotesize\textit{Table \thetable{} \textemdash{} continued from previous page}}\\[3pt]
\toprule
\textbf{Parameter} & \textbf{Value} \\
\midrule
\endhead
\endfoot
\bottomrule
\endlastfoot
Task & Object detection \\
Base model & \texttt{yolov8n.pt} \\
Epochs & 100 \\
Image size & 640 \\
Validation & Enabled \\
Optimizer & Auto-selected \\
Initial learning rate & 0.01 \\
Final learning-rate factor & 0.01 \\
Momentum & 0.937 \\
Weight decay & 0.0005 \\
Warmup epochs & 3.0 \\
Box loss gain & 7.5 \\
Classification loss gain & 0.5 \\
DFL gain & 1.5 \\
Mosaic augmentation & 1.0 \\
Horizontal flip probability & 0.5 \\
Close mosaic & Last 10 epochs \\
\end{xltabular}
}

\subsubsection{Validation and Threshold Selection}

Validation reports precision, recall, F1-score, mAP@50, and mAP@50-95. mAP@50 evaluates
detections at IoU 0.50; mAP@50-95 applies the COCO range from 0.50 to 0.95. A confidence
sweep from 0.05 to 0.95 selects the deployment threshold: lower values preserve more
detections but admit weaker predictions; higher values suppress weak detections but miss more
objects. The operating point is a precision--recall tradeoff, not a single-metric maximum.

\subsubsection{Export and Mobile Post-Processing}

After training, the detector exports to FP32, FP16, and INT8 TFLite variants with the same
class taxonomy. On mobile, the selected variant takes a 640-pixel RGB input. For each
detection row, the post-processor picks the highest class score, multiplies it by objectness
to get confidence, drops candidates below the confidence threshold, and runs non-maximum
suppression at the configured IoU threshold. Surviving boxes are returned with class id, class
name, confidence, and bounding-box coordinates.

\ChapterSubsection{VLM Fine-Tuning and Prediction-Parsing Procedure}

\subsubsection{Adapter Fine-Tuning}

Curated crop-disease records are converted into Qwen-compatible prompt-response examples and
fine-tuned on \texttt{Qwen/Qwen2.5-VL-3B-Instruct} with LoRA adapters
\parencite{hu2021lora}. Adapters target the attention and feed-forward projection layers:
\texttt{q\_proj}, \texttt{k\_proj}, \texttt{v\_proj}, \texttt{o\_proj},
\texttt{gate\_proj}, \texttt{up\_proj}, and \texttt{down\_proj}. Modules matching
\texttt{visual}, \texttt{vision}, or \texttt{merger} are excluded so that only language-side
weights are adapted.

The model loads under 4-bit NF4 quantization with double quantization, bfloat16 compute dtype,
flash attention, and automatic device mapping. Three epochs reflect the available GPU budget,
not a deliberate early stop --- training consumed 12.56~GB of the 22.03~GB available on the
L4. The Liger kernel is enabled; \texttt{torch.compile} is disabled
(Table~\ref{tab:vlm-finetuning-configuration}).

{\small
\renewcommand{\arraystretch}{1.3}
\begin{xltabular}{\textwidth}{@{}ll@{}}
\caption{VLM fine-tuning configuration}\label{tab:vlm-finetuning-configuration}\\
\toprule
\textbf{Parameter} & \textbf{Value} \\
\midrule
\endfirsthead
\multicolumn{2}{l}{\footnotesize\textit{Table \thetable{} \textemdash{} continued from previous page}}\\[3pt]
\toprule
\textbf{Parameter} & \textbf{Value} \\
\midrule
\endhead
\endfoot
\bottomrule
\endlastfoot
Base model & \texttt{Qwen/Qwen2.5-VL-3B-Instruct} \\
Training method & LoRA adapter fine-tuning \\
Quantization & 4-bit NF4 with double quantization \\
Compute dtype & bfloat16 \\
LoRA rank / alpha & 16 / 32 \\
LoRA dropout & 0.05 \\
Epochs & 3 \\
Learning rate & \(2.0 \times 10^{-5}\) \\
Train batch / accumulation & 4 / 4 \\
Effective batch & 16 samples \\
Optimizer & \texttt{paged\_adamw\_8bit} \\
Scheduler & Cosine \\
Max sequence length & 4096 \\
Attention backend & \texttt{flash\_attention\_2} \\
GPU & NVIDIA L4 \\
\end{xltabular}
}

\subsubsection{Prediction Parsing}

At inference time, the gateway receives a crop image and a diagnosis prompt. The generated text
is parsed as a complete JSON object; if that fails, the parser extracts the first embedded JSON
object. A \texttt{canonical\_label} field, if present, becomes the predicted label. Without one,
the prediction is tagged \texttt{\_\_parse\_error\_\_} --- preventing free-form text from
passing as a valid diagnosis and making parseability part of model evaluation.

\ChapterSubsection{Backend Diagnosis-Orchestration Procedure}

\subsubsection{Request Validation and VLM Call}

\texttt{POST /api/v1/diagnose} accepts a multipart request with the image in the
\texttt{image} field and an optional crop-name hint. Validation checks file existence, upload
size, image headers, MIME consistency, dimensions, and format. Images outside
\(64 \times 64\) to \(4096 \times 4096\) pixels are rejected; the crop-name hint is capped at
80 characters. Accepted formats are JPEG, PNG, and WEBP.

Once validated, the diagnosis service checks VLM gateway readiness at \texttt{/ready}. If the
gateway is unconfigured, unreachable, timed out, or not ready, the backend raises a
service-unavailable error instead of returning a partial result. Otherwise it sends an
authorized \texttt{/v1/inference} request with base64 image data, media type, request id,
deterministic generation settings, a 320-token output limit, and a crop-disease prompt.
Authentication uses either a static bearer token or a cached RS256 JWT.

\subsubsection{Advisory Generation and Response Assembly}

The backend parses the VLM prediction with the JSON rules above, then forwards the parsed (or
raw) context and validated image to the Google AI advisory client. The advisory response must
pass schema validation before acceptance. Confidence is clamped to \([0, 1]\), alternatives are
capped at three, and a confidence tier is assigned: \(\geq 0.85\) maps to high, 0.60--0.85 to
medium, lower non-OOD values to low, and out-of-distribution or below 0.30 to OOD. The
assembled response carries request metadata, disease fields, symptoms, advisory content,
processing time, and model version.

\ChapterSubsection{Integrated Online Diagnosis Procedure}

Table~\ref{tab:integration-control-flow} traces the full request lifecycle from image capture
to farmer-facing report.

{\small
\renewcommand{\arraystretch}{1.3}
\begin{xltabular}{\textwidth}{@{} l L L L @{}}
\caption{Integrated online diagnosis control flow}\label{tab:integration-control-flow}\\
\toprule
\textbf{Step} & \textbf{Source component} & \textbf{Implementation evidence} & \textbf{Output} \\
\midrule
\endfirsthead
\multicolumn{4}{l}{\footnotesize\textit{Table \thetable{} \textemdash{} continued from previous page}}\\[3pt]
\toprule
\textbf{Step} & \textbf{Source component} & \textbf{Implementation evidence} & \textbf{Output} \\
\midrule
\endhead
\endfoot
\bottomrule
\endlastfoot
1 & Mobile scan workflow & Capture, preview, editor, analysis view models & Image URI and optional crop context \\
2 & Mobile diagnosis hook & Network-aware mutation with online, slow, offline branches & Cloud request, cached result, or queued scan \\
3 & Mobile API client & Multipart \texttt{/api/v1/diagnose} request with retry and timeout & Backend diagnosis request \\
4 & Backend diagnosis route & Single-image upload middleware and diagnosis controller & Validated service input \\
5 & Backend validation & MIME, header, size, dimension, and crop-name checks & Validated image object \\
6 & VLM client & \texttt{/ready} then \texttt{/v1/inference} with prompt and bearer/JWT auth & Raw VLM prediction text and metadata \\
7 & Backend parser & JSON extraction with raw text preservation & Parsed VLM prediction context \\
8 & Advisory client & Google AI request with strict schema validation & Structured diagnosis and advisory payload \\
9 & Backend response assembly & Request ID, timing, confidence tier, model version & Application-facing JSON response \\
10 & Mobile result workflow & Result view model, history, reminders, feedback, PDF export & Farmer-facing report and local record \\
\end{xltabular}
}

\ChapterSubsection{Mobile Offline and Result-Handling Procedure}

\subsubsection{Scan State and Network Branching}

After image capture (or gallery selection), preview, optional editing, and quality review, the
analysis view model enters one of seven phases: \texttt{ready}, \texttt{scanning},
\texttt{reviewed}, \texttt{sending}, \texttt{success}, \texttt{failed}, or \texttt{queued}.
Network state determines the next step: online sends a cloud request; slow races the request
against a five-second timeout and falls back on failure; offline skips the backend and queues
the scan.

Queued scans are written to SQLite scan history and the unified sync queue as
\texttt{scan\_diagnosis} entries. A cached advisory for the crop context, if available, can
fill in while the scan waits for cloud diagnosis. The API client applies a 60-second timeout
and retries transient failures (408, 429, 500, 502, 503, 504) with exponential backoff.

\subsubsection{Farmer-Facing Result Mapping}

The result view maps completed responses into bilingual disease details, confidence display,
treatment timelines, prevention advice, safety warnings, reminders, feedback, PDF export, and
share cards. Labels appear in Nepali or English based on the current language preference,
treatment steps render as timeline items with optional organic alternatives, and the completed
scan is saved for later access from crop or history screens.
